\chapter{Conclusion}
\label{ch:conclusion}

\section{Summary}

This thesis measured what is lost when the missing entries of a panel of daily
equity returns are filled in. Fifty S\&P~500 constituents over sixteen years
were corrupted under three mechanisms at five severities each, reconstructed by
seven methods, and scored both against the truth directly and through two models
a user would build on the result. We return to the four questions of
Section~\ref{sec:questions}.

The mechanism matters more than the amount. Under the two random patterns the
methods hold up even when a third of the panel is gone: at sparse rate $0.3$ the
market factor is still recovered to within a degree and a half of the true one.
Under the stressed pattern, which removes turbulent dates, the same methods
break down at far smaller shares of missing data. At the heaviest level none
brings the relative covariance error below $0.6$, the first principal component
falls from $41.74\%$ to about $25\%$, and the realised portfolios end at little
more than half of the oracle's cumulative return, with the five imputers
indistinguishable from one another. What has been deleted there is not a random
part of the sample but the part that carries the covariance structure, and no
model can recover from the observed panel what the observed panel does not
contain. A comparable share of missing cells is not a comparable problem.

The deep models do not earn their cost. EM is rarely beaten in the metric tables
and never by a wide margin, CSDI matches it, and BRITS falls behind at two of
the stressed levels. EM fills a panel in about a second, while CSDI needs some
twenty-four minutes of sampling after four hours of training, four orders of
magnitude more for a reconstruction that is not better. The comparison is tied
to the training budget and the hardware of Section~\ref{sec:costs}, but the size
of the gap leaves a good deal of room before the conclusion would turn.

Pointwise accuracy is a poor proxy for the quality of an imputed panel. Under
the MCAR patterns the MAE of mean replacement is only about fifty percent worse
than the best, while its covariance error is close to an order of magnitude
worse than EM's, so a method can look acceptable entry by entry and still
distort the geometry of the panel. The distortion is also systematic rather than
noisy. Every conditional imputer raises the share of variance carried by the
first factor, from $41.74\%$ to between $45.6\%$ and $48.5\%$ at sparse rate
$0.3$, because the value it inserts is a conditional mean and therefore lies
closer to the dominant directions of the panel than a real return does. The more
entries are filled, the more the panel comes to look like its own factor model,
and no pointwise metric registers this.

The rankings largely survive downstream, and the exceptions carry the most
information. Low-rank completion is among the most accurate methods entry by
entry and has the largest five-factor subspace distance at ten of the fifteen
levels, reaching $0.39$ at sparse rate $0.3$ against $0.21$ for EM, since
truncating the spectrum discards exactly the weakly separated middle directions.
Complete-case deletion reads as merely inefficient in the metric tables and
produces allocations whose Sharpe ratio falls to almost zero at sparse rate
$0.1$ and whose mean return turns negative at the gap level $(1000,\,20)$. The
two downstream evaluations are themselves connected: since
$\hat{\bSigma}^{-1} = \sum_i \hat\lambda_i^{-1} \hat v_i \hat v_i^\top$, the
portfolio weights each direction by the reciprocal of its variance and so loads
most heavily on the trailing eigenvectors, which are the directions the
imputations recover worst. The error maximisation seen in
Section~\ref{sec:results_portfolio} is that distortion viewed through an inverse
that magnifies it.

The practical reading is that an imputation should be scored against the use it
is put to. Reconstruction error orders the methods correctly at the easy levels,
where the ordering hardly matters, and is least informative exactly where the
choice of method is consequential.

\section{Limitations}

The evidence comes from a single panel: fifty liquid large-cap equities from one
market over one period, at daily frequency. Fifty assets against $4125$ dates is
a comfortable ratio for covariance estimation, and a cross section closer in size
to the sample length would make every estimate less stable and would likely widen
the differences reported here. The missingness is injected rather than observed,
so its mechanism is known and stationary, whereas real gaps arise from halts,
delistings and illiquidity and follow no law the experimenter can write down; the
stressed pattern is a stylised form of MNAR and not a model of any particular
cause. Each configuration was run once, so the tables carry no estimate of
sampling variation around the numbers they report. The two downstream evaluations
both read the covariance matrix, which makes them complementary but not
independent tests. Finally, all seven methods are smoothers. BRITS reads the
panel in both directions and CSDI conditions on every observed value, so the
results describe reconstruction after the fact and not what would be available to
a user filling a gap in real time, with only the past in hand.

\section{Future Work}

The last of those limitations is the natural continuation. Restricting each
imputer to the information available up to the date being filled turns the
problem from smoothing into a filtering one and makes the cost of that
restriction measurable, which is what a user working in real time needs to know
and what the benchmarks in the literature do not report. A second direction is
to widen the downstream evaluations beyond the second moment, with a
Value-at-Risk backtest, a hedging exercise or a trading rule, since a criterion
that does not pass through the covariance would test parts of the reconstruction
that neither of the present evaluations reaches. A third is to carry the
uncertainty of the imputation forward instead of discarding it. CSDI produces a
predictive distribution that we collapsed to a median, and the inflation of the
factor structure documented above is precisely the artefact that averaging
conditional means creates, so a multiple-imputation treatment in which the
downstream decision sees the spread rather than a point estimate should correct
it. Beyond these, the same design applies directly to larger cross sections,
lower frequencies and other asset classes, and eventually to panels whose gaps
are real rather than injected, where the mechanism has to be inferred rather
than assumed.
